%%%%%%%%%%  scan idx 40  |  folio 27  %%%%%%%%%%
soient $(\rho_i)_{1\leq i\leq r'}$ les classes de représentation de $G$ correspondant aux représentations \emph{fondamentales} de $G/\rad G$, et $(\sigma_i)_{1\leq i\leq r''}$ une base du groupe des homomorphismes rationnels de $G$ dans $k^*$~; on a donc $r'+r''=r$ (rang de $G$).

\textbf{Corollaire :} \emph{L'anneau $\underline{K}_r(G)$ s'identifie au sous-anneau}
\[
\Z[\rho_1,\ldots,\rho_{r'};\sigma_1,\ldots,\sigma_{r''};\sigma_1^{-1},\ldots,\sigma_{r''}^{-1}]
\]
\emph{du corps des fonctions rationnelles en les $\rho_i$ et les $\sigma_j$ à coefficients dans $\Q$}\footnote{Il faut supposer le groupe dévisé $G'$ de $G$ simplement connexe pour que cet énoncé soit correct.}.

En particulier, $\underline{K}_r(G)$ n'a pas de diviseurs de $0$~; si $G=\prod_i\Gl(n_i,k)$, l'anneau $\underline{K}_r(G)$ s'explicite ainsi~: soit $\rho_i$ la représentation de $G$ déduite de la représentation identique de $\Gl(n_i,k)$~; si l'on considère alors le corps des fonctions rationnelles sur $\Q$ en des indéterminées $\lambda^j(\rho_i)$ ($1\leq j\leq n_i$), alors $\underline{K}_r(G)$ s'identifie au sous-anneau engendré par ces indéterminées et les inverses des $\lambda^{n_i}(\rho_i)$. Cela justifie notre assertion que, du moins en caractéristique $0$, les opérations $\lambda^i$ permettent de reconstituer toutes les opérations tensorielles sur les fibrés vectoriels~; ceci est encore vrai sur un corps algébriquement clos de caractéristique quelconque, pourvu qu'on prenne des classes réduites.

\textbf{Remarques :} a) Dans l'exemple 1), l'anneau $\underline{K}(G)$ est augmenté, l'augmentation étant définie par le degré des représentations, et par passage au quotient, cette augmentation définit une augmentation sur $\underline{K}_r(G)$. De même, les $\lambda$-anneaux envisagés dans l'exemple 2) sont augmentés lorsque $X$ est $k$-irréductible~; dans $\underline{K}(\xi(X))$, l'augmentation est définie par le \emph{rang d'un fibré vectoriel}, c'est-à-dire la dimension de ses fibres.

